\SourceLeaf{355}%
\par\indent When a triangulation network on the spheroid, whose sides are so-called geodesic lines, is mapped conformally onto the sphere, it is represented by a corresponding triangulation network. As follows directly from the meaning of conformality, its angles are \textit{exactly} equal to the corresponding angles in the original network. Its sides, though very close to great-circle arcs, differ from them in the strict mathematical sense. Suppose these deviations can be made negligible throughout the entire network at the accuracy required for the calculation. Then, once one side of the spheroidal network has been mapped onto the sphere, the whole network can clearly be calculated directly as an ordinary spherical triangulation. Only at the end is it necessary to convert from longitudes and latitudes on the sphere back to longitudes and latitudes on the spheroid, if the final results of the measurements are wanted in that form. This conversion can be made either with the formulas provided by the chosen mapping or with a table calculated in advance. If, on the other hand, the deviation is large enough to be worth taking into account, each angle obtained from the measurements will first need a small reduction before the precise calculation on the sphere. This will add little to the work if the numerical values of the reductions can be calculated easily.

With the mapping developed in this memoir, the deviation of the curve that represents a geodesic arc on the sphere from the great-circle arc between the same endpoints can always be treated as negligible near a certain parallel (the standard parallel). This parallel can be chosen at will. If one sets up the entire calculation oneself from the beginning for a given triangulation network, the most suitable choice may be to put the parallel roughly through the middle of the whole network. The farther one goes north or south from this standard parallel. the greater these deviations may become. which, in addition, depend both on the lengths of the triangle sides and on their orientation relative to the meridian. Even at a very considerable distance from the standard parallel, however, they are always still so %
